\documentclass[10pt,a4paper]{amsart}
\usepackage[margin=19mm]{geometry}
\usepackage{amsmath,amssymb,mathtools}
\usepackage[scr=rsfs,cal=euler]{mathalfa}
\usepackage{xfrac}
\usepackage[unicode,hidelinks,bookmarks=false]{hyperref}
\newcommand{\ie}{i.e.\ }
\newcommand{\cf}{cf.\ }
\newcommand{\resp}{respectively\ }
\newcommand{\Subsection}{\subsection}
\providecommand{\nobreakdash}{\nobreak\hbox{-}}
\setlength{\emergencystretch}{2em}
\multlinegap0pt
\usepackage{amssymb}
\usepackage{xcolor}
\newtheorem{lem}{Lemma}
\newtheorem{prop}{Proposition}
\def\R{\mathbb{R}}
\def\d{\mathrm{d}}
\title{Eigenforms and special holonomy}
\author{Bobby Acharya, Luc\'ia Cabrera, Simone Corbo, Jason D. Lotay}
\date{Source: arXiv:2608.09878v2 (CC BY 4.0)}
\begin{document}
\maketitle

\noindent\emph{Source and licence.} Eigenforms and special holonomy. Version arXiv:2608.09878v2 (CC BY 4.0), distributed by arXiv under the Creative Commons Attribution 4.0 International licence, http://creativecommons.org/licenses/by/4.0/, as recorded in the arXiv licence metadata for this exact version and shown on the abstract page https://arxiv.org/abs/2608.09878v2. Full licence notice: \texttt{licenses/arxiv-2608.09878-CC-BY-4.0.txt}.\par\smallskip

\noindent\textit{Editorial scope.} Counted unit: the article's prop prop:hk.ALE.1forms (source0.tex lines 680--681). The article's complete original proof of it is delivered with it (source0.tex lines 683--706, one \texttt{proof} environment of 24 lines). No mathematics is written, repaired or re-derived: the statement and the proof are the article's own lines, in the article's own notation, and no retained line is substituted or re-flowed.  Retained uncounted as the source-local context the proof needs: lines 254--256; lines 264--266; lines 397--399; lines 472--474; lines 622--624; lines 645--647.  Nothing was omitted from any retained span, so the package carries no omission record.  No retained line contains a citation, so the wrapper carries no bibliography and no bibliography entry.  No retained line contains a figure environment, an image-inclusion command or drawing code, so no image is delivered and the package declares no asset dependency.  Editorial formatting only, in addition: an AMS \texttt{amsart} preamble that restates the article's own declarations and macros used by the retained lines, namely the environments \texttt{lem}, \texttt{prop} on separate counters, together with the standard packages those lines need.  The retained environments print in this document as Lemma 1, Proposition 1 in order of appearance, whereas the article prints the same items with the numbering of its own full document.  The complete native source of the article as distributed in the arXiv e-print for this exact version is bundled unchanged as \texttt{sources/207280/source0.tex}.
\begin{equation}\label{eq:eigenform}
    \Delta\alpha=\lambda\alpha
\end{equation}

\begin{equation}\label{eq:normalizability.1}
    \alpha\in L^2\Omega^k(M,g),
\end{equation}

\begin{equation}\label{eq:lambda.non.neg}
\lambda\|\alpha\|_{L^2}^2=\langle \alpha,\Delta\alpha\rangle_{L^2}=\|\d\alpha\|_{L^2}^2+\|\d^*\alpha\|_{L^2}^2\geq 0,
\end{equation}

\begin{lem}\label{lem:integrability.grav.inst}  Let $(M^4,g)$ be a gravitational instanton.  
   Let $\alpha\neq 0$ be a differential form satisfying \eqref{eq:eigenform} for $\lambda\in\R$ and \eqref{eq:normalizability.1}.  Then $\d\alpha,\d^*\alpha\in L^2$ and \eqref{eq:lambda.non.neg} holds, so $\lambda\geq 0$.
\end{lem}

\begin{lem}\label{lem:hk.ALE.sd}   
If $(M^4,g)$ is complete hyperk\"ahler then $\mathcal{E}^2_+(\lambda)\neq 0$ if and only if $\mathcal{E}^0(\lambda)\neq 0$.
\end{lem}

\begin{lem}\label{lem:hk.asd.exact}  Let $(M^4,g)$ be a gravitational instanton.  
If $\alpha\in L^2\Omega^1(M)$ and $\d\alpha\in L^2\Omega^2_{\pm}(M)$ then $\d\alpha=0$.
\end{lem}

\begin{prop}\label{prop:hk.ALE.1forms} Let $(M^4,g)$ be a gravitational instanton.  If $\mathcal{E}^0(\lambda)=0$ for $\lambda\neq 0$ then $\mathcal{E}^1(\lambda)=0$.
\end{prop}

\begin{proof}  Let $\alpha\in \mathcal{E}^1(\lambda)$.  
 We first note that by Lemma \ref{lem:integrability.grav.inst} we have $\d^*\alpha,\d\alpha\in L^2$.  We then see that
 \begin{equation}\label{eq:Delta.d*}
      \Delta \d^*\alpha=\d^*\Delta\alpha=\lambda\d^*\alpha.
 \end{equation}
 Hence $\d^*\alpha\in \mathcal{E}^0(\lambda)=0$ by hypothesis.
 
 
 We also see that 
\begin{equation}
     \d_+\alpha=\tfrac{1}{2}(\d\alpha+*\d\alpha)\in L^2\Omega^2_+(M).
\end{equation}
 Again we have that
\begin{equation}
    \Delta \d_+\alpha=\d_+\Delta\alpha=\lambda\d_+\alpha
\end{equation}
and so, by Lemma \ref{lem:hk.ALE.sd}, we know that $\d_+\alpha=0$. We then observe that
\begin{equation}
    \d\alpha=\tfrac{1}{2}(\d\alpha-*\d\alpha)=\d_-\alpha\in L^2\Omega^2_-(M).
\end{equation}
We may deduce that $\d\alpha=0$ from Lemma \ref{lem:hk.asd.exact}.

Given that $\d\alpha=0=\d^*\alpha$, we have that $\Delta\alpha=0$, but as $\lambda\neq 0$ this forces $\alpha=0$ as claimed.
\end{proof}

\end{document}
